\chapter{Conclusion}

At first glance, it seems quite odd that category theory, an abstract branch of mathematics is so closely tied with any programming language. Category theory is all about abstracting patterns and structures that can be used to generalize concepts. Computer science is all about computation and producing results and systems. However, as we've discovered, the two concepts are actually quite closely tied. In fact, if we dig deeper we can see that the nature of computer science is in fact abstraction. Computers processing comes from computer chips, which consist of billions of tiny transistors (electrically controlled switches), that produce basic logic. From thos basic building blocks, we abstract further and further up the stack. It is the programmers jobs to handle these abstractions, building more and more complex structures on top of these basic building blocks. Haskell, as a high-level functional programming language is an amazing environment for these categorical concepts.

First, we saw how the concept of functors can be translated naturally into operations in Haskell. We saw examples of \lstinline{Maybe} and \lstinline{List}. Then, we explored the concept of monoids in Haskell. Next, we saw how monoids are generalized to categories. In fact, monads are monoids in the category of of endofunctors. Finally, we saw the application of monads in Haskell and their general usefulness in computer science. Monads are a powerful tool because they allow us to generalize the idea of composition, the essence of functional programming. Thanks for joining us on this journey!